% !TEX root = ../../tree_notes.tex
% =====================================================================
%  Chapter 9 opener -- "Chapter at a Glance", three reading paths + the
%  ladder staircase, the two running examples (the array A = 3,1,4,2 and
%  the histogram H = 2,1,5,6,2,3), the "who owns what" comparison
%  (monotonic stack vs segment tree vs treap), roadmap, recognition guide,
%  constraint->tool table, notation.
%  NO numbered boxes here. Exports fig:c26-hist and fig:c26-cartA.
% =====================================================================

\begin{chaptertheorybox}{Chapter at a Glance}
Chapter~8 (\chref{c06}) made a tree short \emph{by editing it}: you insert or delete, then pay a
rotation to repair the height. This chapter never edits anything. It starts from a bare
\term{array} and asks a different question -- \emph{what tree is already hiding inside these
numbers?} Hand it an array $a$ and there is exactly \emph{one} binary tree whose inorder reading is
the array in order and whose heap rule is the values: the \term{Cartesian tree}. And there is one
tiny machine that reads that tree straight out of the array in a single left-to-right pass -- a
\term{monotonic stack}. The astonishing part is that this \emph{same} stack, with no tree built at
all, is the whole solution to largest rectangle in a histogram, sum of subarray minimums, next
greater element, sliding-window maximum, range minimum query, and the chapter's Summit. One
invariant, one stack, an entire problem family:
\begin{itemize}[nosep]
  \item \textbf{Two orders, one tree.} A Cartesian tree obeys \emph{BST order on the index} (inorder
        $=$ the array) and \emph{heap order on the value} at the same time. Under distinct values the
        tree is unique, and its root is just the position of the minimum. \hfill \lv{3}
  \item \textbf{The stack is the build.} Sweep the array, keep the stack values increasing, and pop
        while the new element breaks monotonicity -- the popped links \emph{are} the Cartesian tree's
        parent pointers, produced in $\bigO{n}$. The moment a bar stops being the smallest is exactly
        the moment its rectangle or span is decided. \hfill \lv{3}--\lv{5}
  \item \textbf{Span $=$ subtree.} The subtree of node $i$ is precisely the maximal window in which
        $a[i]$ is the minimum. That one lemma converts a hundred "sum of $a[i]\cdot(\text{how many
        ranges where it wins})$'' problems into two stack passes for the prev/next smaller (or
        greater) bounds. \hfill \lv{4}--\lv{5}
\end{itemize}

\textbf{You will learn to}
\begin{itemize}[nosep]
  \item build a Cartesian tree in $\bigO{n}$ with a stack and read the two invariants off it; \hfill\lv{3}
  \item recognise "the answer is decided when the stack pops" and turn histogram, subarray-minimum
        and rain-water problems into one or two stack passes; \hfill\lv{3}--\lv{5}
  \item keep sliding-window extremes with a monotonic \emph{deque}; \hfill\lv{4}
  \item reduce a static \term{range minimum query} to an LCA on the Cartesian tree, and know exactly
        what you are handing off to the LCA machinery of Chapter~16 (\chref{c11}); \hfill\lv{3}--\lv{4}
  \item see a Cartesian tree as a \emph{treap whose priorities are the values} and a random-permutation
        BST as a random Cartesian tree -- borrowing Chapter~8's height bound, not re-proving it. \hfill\lv{4}
\end{itemize}

\textbf{Prerequisites.} Arrays, and the \term{stack} as a container, from Chapter~1--2; the BST
inorder fact from Chapter~7 (\chref{c05}) is \emph{leaned on}, never re-derived. Chapter~8's treap
(\secref{c06-treaps}) and its expected-height theorem (\thmref{c06-exptheight}) are used as a black
box only in \secref{c26-treapview}. LCA is quoted, not built (Chapter~16, \chref{c11}).
\end{chaptertheorybox}

% ---------------- Reading paths + ladder position ----------------
\begin{chaptertheorybox}{How to Read This Chapter (three paths)}
Level badges $\lv1$--$\lv6$: $\lv3$ is a solid LC-medium stack problem, $\lv4$ is a Hard you can only
crack once the "the span is the subtree" idea clicks, and $\lv5$ is an LC/CF Hard where the whole
difficulty is \emph{realising} the monotonic structure is there at all.

\textbf{Fast path ($\lv3$, ``I just want the stack'').} Read \secref{c26-monostack} directly -- the
monotonic stack is a standalone technique you can use without ever naming a Cartesian tree. Do the
solved histogram (LC 84) and the next-greater / sliding-window practices. Skip the tree build, the
RMQ reduction and the treap view. You leave able to cold-solve largest-rectangle, daily-temperatures
and sliding-window-maximum.

\textbf{Structural path ($\lv3$--$\lv4$).} Start at \secref{c26-cartesian} and \emph{build} the tree
-- it makes \secref{c26-monostack} feel inevitable rather than clever, and pays off in
\secref{c26-rmq} where the tree is the point. Recommended order: 9.1 $\to$ 9.2 $\to$ 9.3.

\textbf{Full path.} Everything in order, finishing at \secref{c26-treapview}, where the Cartesian tree
and the Chapter~8 treap turn out to be the \emph{same object}, and at the Summit (CF 817D), which is
the purest expression of the span-equals-contribution idea: the answer is a signed sum of
$a[i]\cdot(\text{span where it is max}) - a[i]\cdot(\text{span where it is min})$, both from a stack
pass each.

The forward rope is Chapter~10 (\chref{c07}) for genuine \emph{heaps} (a priority queue you edit, not
a heap-order you read off a static array), and Chapter~16 (\chref{c11}) for the LCA engine the RMQ
reduction calls.
\end{chaptertheorybox}

\begin{figure}[H]
\centering
\begin{tikzpicture}[diagram, x=9mm, y=7mm]
  \foreach \st/\lab in {1/lv1,2/lv2,3/lv3,4/lv4,5/lv5,6/lv6} {
    \fill[\lab!22, draw=\lab, line width=0.8pt] (\st-1,0) rectangle (\st,\st*0.52);
    \node[font=\scriptsize\bfseries, text=\lab] at (\st-0.5,\st*0.52-0.42) {L\st};}
  \node[annot, anchor=south west, font=\scriptsize]
    at (0.05,0.62) {read an array as a tree};
  \node[annot, anchor=south west, font=\scriptsize]
    at (1.05,1.14) {9.1 build; 9.2 stack basics};
  \node[annot, anchor=south west, font=\scriptsize]
    at (2.05,1.66) {9.2 histogram, next-greater, deque};
  \node[annot, anchor=south west, font=\scriptsize]
    at (3.05,2.18) {9.1/9.3 span=subtree, RMQ};
  \node[annot, anchor=south west, font=\scriptsize, text=lv5]
    at (4.05,2.70) {9.2 LC 84; 9.3 RMQ lab; treap view};
  \node[annot, anchor=south west, font=\scriptsize, text=lv5]
    at (5.05,3.22) {summit (CF 817D)};
\end{tikzpicture}
\caption*{Chapter~9 on the six-step ladder: seeing the array$\to$tree reading is $\lv3$; the monotonic
stack and histogram sit at $\lv3$--$\lv4$; the ``span $=$ subtree'' contribution trick and the RMQ
reduction climb to $\lv4$; solving LC 84 and the treap-equivalence view reach $\lv5$, where the Summit
(CF 817D \emph{Imbalanced Array}) turns the whole chapter into two signed stack passes on $n\le 10^6$.
Fast path $=$ steps 2--3; full path runs to 6.}
\end{figure}

% ---------------- The running examples ----------------
\subsection*{The Running Examples: One Array, One Histogram}

Two small inputs, reused in every section so the ideas always land on something concrete.

\textbf{Array $A$} $= (3, 1, 4, 2)$, positions $1,2,3,4$. Its values are \emph{distinct}, so its
min-Cartesian tree is unique: the smallest value, $1$ at position $2$, is the root; the elements to
its left form its left subtree and those to its right form its right subtree, recursively. Read the
tree's inorder and you get $A$ back; read any node against its children and the parent is smaller --
both orders at once. $A$ drives the build (\secref{c26-cartesian}), the next-smaller spans, the RMQ
queries (\secref{c26-rmq}) and the treap equivalence (\secref{c26-treapview}).

\textbf{Histogram $H$} $= (2, 1, 5, 6, 2, 3)$ -- the official LeetCode~84 input, whose answer is the
rectangle of height $5$ and width $2$ over bars $3,4$, giving area $\boxed{10}$. $H$ has a
\emph{duplicate} ($2$ appears twice), so it also shows why tie-breaking must be asymmetric when spans
are counted. It drives the monotonic-stack section (\secref{c26-monostack}) and reappears in the
Summit's contribution argument.

\begin{figure}[H]
\centering
\begin{tikzpicture}[diagram, x=11mm, y=6mm]
  % histogram H = 2,1,5,6,2,3
  \def\hs{{2,1,5,6,2,3}}
  \draw[thick, darkgray] (0,0) -- (6.6,0);
  \foreach \i/\h in {0/2,1/1,2/5,3/6,4/2,5/3} {
    \fill[primaryblue!22, draw=primaryblue, thick] (\i,0) rectangle (\i+1,\h);
    \node[font=\scriptsize, text=darkgray] at (\i+0.5,-0.45) {\i+1};
    \node[font=\scriptsize\bfseries, text=deepblue] at (\i+0.5,\h+0.35) {\h};}
  % the winning rectangle height 5 over bars 3,4 (x=2..4)
  \draw[ultra thick, warningred, dashed] (2,0) rectangle (4,5);
  \node[font=\scriptsize\bfseries, text=warningred, align=center] at (3,2.6) {$5\times 2$\\ area $10$};
  \node[annot, align=left, font=\scriptsize] at (7.0,3) {\textbf{Histogram $H$}\\ $2,1,5,6,2,3$\\
    the tallest rectangle is\\ the dashed one, area $10$};
\end{tikzpicture}
\caption{The histogram running example $H=(2,1,5,6,2,3)$. Bars are drawn in index order; the
monotonic stack (\secref{c26-monostack}) finds that the widest rectangle of height $5$ spans bars
$3$--$4$ and has area $10$ -- the LeetCode~84 answer.}\label{fig:c26-hist}
\end{figure}

\begin{figure}[H]
\centering
\begin{tikzpicture}[diagram, scale=0.95, transform shape]
  % min-Cartesian tree of A = 3,1,4,2 (index shown as subscript)
  \node[troot] (a2) at (0,0)   {$a_2{=}1$};
  \node[tleaf] (a1) at (-2,-1.4) {$a_1{=}3$};
  \node[tnode] (a4) at (2,-1.4)  {$a_4{=}2$};
  \node[tleaf] (a3) at (1,-2.8)  {$a_3{=}4$};
  \draw[tedgehl] (a2) -- (a1);
  \draw[tedgehl] (a2) -- (a4);
  \draw[tedgehl] (a4) -- (a3);
  \node[annot, align=left, font=\scriptsize] at (-4.6,-1.4)
    {\textbf{Inorder} (left,root,right):\\ $a_1,a_2,a_3,a_4$\\ $=3,1,4,2$ -- back to $A$};
  \node[annot, align=left, font=\scriptsize] at (5.0,-2.0)
    {\textbf{Min-heap on value:}\\ parent $<$ both children\\ $1<3,1<2<4$};
\end{tikzpicture}
\caption{The min-Cartesian tree of $A=(3,1,4,2)$. The root is the minimum ($1$, at position $2$); the
left part ($a_1$) and right part ($a_3,a_4$) recurse. Its inorder is the array itself, and every
parent is smaller than its children -- the two invariants of \defref{c26-cartesian}, shown on one
object. Node $a_4$'s subtree is the single window $\{3,4\}$ where $2$ is the minimum: the
span-equals-subtree idea of \lemref{c26-span}.}\label{fig:c26-cartA}
\end{figure}

\medskip
\noindent Three engines answer the \emph{same} range-extremum questions; knowing which owns what stops
you re-solving a solved problem:
\begin{center}
\renewcommand{\arraystretch}{1.25}
\setlength{\tabcolsep}{5pt}
\begin{tabular}{@{}p{0.185\linewidth} p{0.42\linewidth} p{0.31\linewidth}@{}}
\toprule
\textbf{Engine} & \textbf{What it is for} & \textbf{Where it lives} \\
\midrule
Monotonic stack / deque & \emph{offline / one-pass} questions where the answer is decided the moment a
  bar is popped: next greater, histogram spans, sliding-window extreme & this chapter (\secref{c26-monostack}) \\
Segment / sparse table  & \emph{arbitrary} online range queries (min, sum, \dots) over a
  \emph{changing} array & Chapter~18--19 (\chref{c14}, \chref{c13}) \\
Treap (heap on random priority) & a \emph{dynamic} ordered set you insert into and delete from while
  staying short & Chapter~8 (\chref{c06}) \\
\bottomrule
\end{tabular}
\end{center}

% ---------------- Roadmap diagram ----------------
\begin{figure}[H]
\centering
\begin{tikzpicture}[diagram, node distance=3.5mm,
    top/.style={boxnode, fill=primaryblue!10, draw=primaryblue, text width=19mm,
                minimum height=11mm, font=\small},
    res/.style={annot, text width=19mm, anchor=north, font=\scriptsize}]
  \node[top] (t1) {\textbf{9.1} Cartesian Build};
  \node[top, right=of t1] (t2) {\textbf{9.2} Monotonic Stack};
  \node[top, right=of t2] (t3) {\textbf{9.3} RMQ via LCA};
  \node[top, right=of t3] (t4) {\textbf{9.4} Random Treap View};
  \foreach \a/\b in {t1/t2, t2/t3, t3/t4} \draw[flow] (\a) -- (\b);
  \node[res] at ([yshift=-2mm]t1.south) {array $\to$ tree; the two orders};
  \node[res] at ([yshift=-2mm]t2.south) {stack \emph{is} the build; histogram};
  \node[res] at ([yshift=-2mm]t3.south) {range min $=$ LCA value};
  \node[res] at ([yshift=-2mm]t4.south) {Cartesian $=$ treap; random BST};
\end{tikzpicture}
\caption*{The four topics are one picture seen four ways. 9.1 \emph{builds} the tree and names its two
invariants; 9.2 reveals the build is just a monotonic stack and runs the whole application family on
that stack; 9.3 uses the tree as an \emph{index} and reduces range minimum to LCA; 9.4 steps back and
shows the tree is secretly the Chapter~8 treap with values-as-priorities -- tying the chapter's ending
to where the balanced-tree story closed.}
\end{figure}

% ---------------- Topic table ----------------
\begin{center}
\renewcommand{\arraystretch}{1.25}
\setlength{\tabcolsep}{4pt}
\begin{tabular}{@{}p{0.235\linewidth} p{0.365\linewidth} c p{0.23\linewidth}@{}}
\toprule
\textbf{Topic} & \textbf{Key results} & \textbf{Levels} & \textbf{Solved problem} \\
\midrule
\secref{c26-cartesian}\ \ Cartesian Tree in $O(n)$
  & two orders (inorder=index, heap=value); unique on distinct values; stack build & $\lv3$--\lv{4}
  & \probref{c26-goodsub} \\
\addlinespace
\secref{c26-monostack}\ \ Monotonic Stack Applications
  & pop $=$ decision; histogram, next-greater, deque sliding max, rain water & $\lv3$--\lv{5}
  & \probref{c26-hist84} \\
\addlinespace
\secref{c26-rmq}\ \ Range Minimum Query via the Tree
  & RMQ of $a[l..r]$ $=$ value at $\lca(l,r)$ in the Cartesian tree & $\lv3$--\lv{4}
  & \probref{c26-rmq1647} \\
\addlinespace
\secref{c26-treapview}\ \ The Cartesian Tree as a Treap
  & values $=$ priorities $\Rightarrow$ treap; permutation $\Rightarrow$ random BST & $\lv4$
  & \probref{c26-mbt2} \\
\bottomrule
\end{tabular}
\end{center}

% ---------------- Problem recognition guide ----------------
\subsection*{Problem Recognition Guide --- When You See \emph{this}, Think \emph{that}}

Monotonic-stack problems are the ones where an element's fate is sealed by the first \emph{smaller}
(or \emph{larger}) neighbour to its left or right. If the phrasing mentions a threshold, a window, or
"extends until", that is the tell; every row points at the section that builds the tool.

\begin{center}
\renewcommand{\arraystretch}{1.3}
\setlength{\tabcolsep}{4pt}
\begin{tabular}{@{}p{0.44\linewidth} p{0.34\linewidth} l@{}}
\toprule
\textbf{When the problem says\ldots} & \textbf{\textcolor{deepblue}{\textbf{Think\ldots}}} & \textbf{Section} \\
\midrule
"largest rectangle in a histogram / maximal binary grid"
  & pop when the bar stops being a minimum; width fixed at pop & \secref{c26-monostack} \\
\addlinespace
"sum over all subarrays of the min (or max)" / "how many ranges is $a[i]$ the winner of"
  & prev/next smaller (or greater) span; $a[i]\cdot(i-L)\cdot(R-i)$ & \secref{c26-cartesian}, Summit \\
\addlinespace
"next greater element", "daily temperatures", "can you see a taller building"
  & decreasing stack; pop when a bigger value arrives & \secref{c26-monostack} \\
\addlinespace
"max (or min) inside a \emph{sliding window} of size $k$"
  & monotonic \emph{deque}; drop from both ends & \secref{c26-monostack} \\
\addlinespace
"minimum of a static range $[l,r]$, many queries"
  & value at $\lca(l,r)$ in the Cartesian tree & \secref{c26-rmq} \\
\addlinespace
"rebuild a max-tree / insert at the far right"
  & the Cartesian / treap view; walk the right spine & \secref{c26-treapview} \\
\bottomrule
\end{tabular}
\end{center}

\begin{center}
\renewcommand{\arraystretch}{1.25}
\setlength{\tabcolsep}{4pt}
\begin{tabular}{@{}p{0.29\linewidth} p{0.42\linewidth} p{0.20\linewidth}@{}}
\toprule
\textbf{Constraint you see} & \textbf{What it forces} & \textbf{Chapter-9 tool} \\
\midrule
$n \le 10^{5}$--$10^{6}$, "sum over all subarrays" & $\bigO{n}$ -- an $\bigO{n^2}$ window scan dies & stack span \\
\addlinespace
answer decided by nearest greater / smaller neighbour & compute all bounds in one pass each & \secref{c26-monostack} \\
\addlinespace
window of fixed width $k$, extreme per window & monotonic deque, amortised $\bigO{1}$ per step & \secref{c26-monostack} \\
\addlinespace
static array, many range-min queries, few updates & Cartesian tree $+$ LCA (or sparse table) & \secref{c26-rmq} \\
\addlinespace
products that overflow $32$ bits ($a[i]\cdot\text{counts}$) & accumulate in \code{long long} & Summit (\probref{c26-summit}) \\
\bottomrule
\end{tabular}
\end{center}

\begin{notebox}[Not this chapter?]
If the structure must \emph{change} -- insert, delete, then ask a range question online -- a one-pass
stack cannot remember your edits; that is the segment tree (\chref{c13}) or a balanced tree
(\chref{c06}). If you only need the \emph{global} largest or smallest over the whole array and nothing
per-window, a single variable suffices -- no stack. If you want the sliding-window \emph{median}, that
multiset / two-heap version is Chapter~8 (\chref{c06}); this chapter owns the sliding
\emph{extreme} via the deque. And if the tree must be \emph{balanced against an adversary}, you need
rotations -- Chapter~8 again. This chapter owns \emph{reading a static array as a tree, and the one
stack that does all the work}.
\end{notebox}

% ---------------- Notation ----------------
\subsection*{Notation Used in This Chapter}

\begin{center}
\renewcommand{\arraystretch}{1.2}
\setlength{\tabcolsep}{4pt}
\begin{tabular}{@{}p{0.20\linewidth} p{0.27\linewidth}@{\quad}p{0.17\linewidth} p{0.27\linewidth}@{}}
\toprule
\textbf{Symbol} & \textbf{Meaning} & \textbf{Symbol} & \textbf{Meaning} \\
\midrule
$A$                 & the running array $(3,1,4,2)$ (\figref{c26-cartA}) & $H$ & the running histogram $(2,1,5,6,2,3)$ (\figref{c26-hist}) \\
$a[i]$              & value at index $i$; indices are $1$-based here & $\mathit{stack}$ & the monotonic stack of indices \\
$L_i,\ R_i$         & nearest index to the left/right of $i$ that \emph{beats} $a[i]$ (prev/next smaller or greater) & $\mathit{span}_i$ & $(L_i, R_i)$, the window where $a[i]$ is the extremum \\
$\lca(u,v)$         & lowest common ancestor in a tree (\chref{c11}); here in the Cartesian tree & $T$ & the Cartesian tree of an array \\
$\text{inorder}(T)$ & the array read off the tree by an inorder walk & $\bigstar$ & marks the Summit \\
\bottomrule
\end{tabular}
\end{center}
